% GENERATED FILE. Edit canonical lessons, then run npm run generate:books.
% canonical-source-hash: fnv1a64:4b5e96b195889104
% canonical-lessons: ZH-C471-yanlei, ZH-C471-yanqian, ZH-W471-yan-eye, ZH-C471-xiaoxinyan, ZH-C471-jinshi, ZH-C471-zhayan

\chapter{Eyes}
\label{ch:zh-eyes}
\hlchaptermodality{471}

\begin{chapteropening}
\textbf{By the end of this chapter:} \emph{I can talk about tears, what is right in front of me, petty people, being short-sighted and a wink.}

Complete the last lesson of chapter 471: I can talk about tears, what is right in front of me, petty people, being short-sighted and a wink.
\end{chapteropening}

\section[yǎnlèi]{yǎnlèi — tears}

\label{lesson:ZH-C471-yanlei}



\begin{quote}
Before the new one: say the Mandarin for slapdash, then the Mandarin for slow-witted.
\end{quote}

\subsection*{You'll want to know: yǎnlèi}
\emph{yǎnlèi} — \textquotedblleft{}tears\textquotedblright{}.

\begin{quote}
\emph{wǒ dì-yī cì kàn dào bàba de yǎnlèi} — It was the first time I had seen my dad in tears.
\end{quote}

Its characters come later; for now it is a word for the ear.

\begin{cousinweb}
\emph{yǎnlèi} is built from \emph{yǎn} 'eye' and \emph{lèi} 'tears', both inherited from Middle Chinese, \emph{ngeanX} and \emph{lwijH} (Guangyun, in Baxter's transcription), which glosses the first 'the eye' and the second 'tears'.
\end{cousinweb}

\begin{grammarlens}[title={What you've built}]
One more word for saying how it was.
\end{grammarlens}

\subsection*{Your turn}
\begin{itemize}
\raggedright
  \item \emph{Say it:} \emph{yǎnlèi}
  \item \emph{Say it:} \emph{yǎnlèi}, once more
  \item \emph{Say it:} \emph{wǒ dì-yī cì kàn dào bàba de yǎnlèi} — It was the first time I had seen my dad in tears.
  \item \emph{Recall:} say the Mandarin for slapdash, then the Mandarin for slow-witted, then say \emph{yǎnlèi} again
\end{itemize}

\begin{tcolorbox}[breakable,colback=teal!4,colframe=teal!35!black,title={Before you move on}]
What does yǎnlèi mean? (\textquotedblleft{}Tears\textquotedblright{}.) Say it once more.
\end{tcolorbox}
\section[\texorpdfstring{\zh{眼前} (yǎnqián)}{眼前 (yǎnqián)}]{\zh{眼前} (yǎnqián) — before one's eyes; right now}

\label{lesson:ZH-C471-yanqian}



\begin{quote}
Before the new one: say the Mandarin for slow-witted, then the Mandarin for tears.
\end{quote}

\subsection*{You'll want to know: \zh{眼前}}
\textbf{\zh{眼前}} — \emph{yǎnqián} — \textquotedblleft{}before one's eyes; right now\textquotedblright{}.

\begin{quote}
\textbf{\zh{那个事故发生在我眼前}} \emph{nàge shìgù fāshēng zài wǒ yǎnqián} — That accident happened right in front of my eyes.
\end{quote}

The new character, in pieces: \zh{眼} — eleven strokes in two pieces: on the left an eye, five strokes, like \zh{日} \emph{rì} with one bar more; on the right six strokes, a stroke that runs right and turns down, two bars, a vertical that turns up at the foot, a short stroke falling to the left and a long one falling to the right. The next lesson puts it in your hand.

\begin{cousinweb}
\textbf{\zh{眼前}} (\emph{yǎnqián}) is built from \zh{眼} 'eye' and \zh{前} 'in front', both inherited from Middle Chinese, \emph{ngeanX} and \emph{dzen} (Guangyun, in Baxter's transcription), which glosses the first 'the eye' and the second 'before'.
\end{cousinweb}

\begin{grammarlens}[title={What you've built}]
One more word for saying how it was.
\end{grammarlens}

\subsection*{Your turn}
\begin{itemize}
\raggedright
  \item \emph{Say it:} \emph{yǎnqián}
  \item \emph{Say it:} \emph{yǎnqián}, once more
  \item \emph{Say it:} \emph{nàge shìgù fāshēng zài wǒ yǎnqián} — That accident happened right in front of my eyes.
  \item \emph{Recall:} say the Mandarin for slow-witted, then the Mandarin for tears, then say \emph{yǎnqián} again
\end{itemize}

\begin{tcolorbox}[breakable,colback=teal!4,colframe=teal!35!black,title={Before you move on}]
What does \zh{眼前} mean? (\textquotedblleft{}Before one's eyes; right now\textquotedblright{}.) Say it once more.
\end{tcolorbox}
\section[\texorpdfstring{\zh{眼} (yǎn)}{眼 (yǎn)}]{\zh{眼} — an eye, then a right side of six}

\label{lesson:ZH-W471-yan-eye}



\begin{quote}
The word first, then the character inside it.

\begin{itemize}
\raggedright
  \item \emph{Recall:} say \emph{yǎnqián} — \textbf{R1}, one lesson back, before it decays
\end{itemize}

\begin{quote}
\zh{日} \zh{前}
\end{quote}

\emph{Read it:} each one aloud, with nothing said first — \emph{rì}, \emph{qián}
\end{quote}

\subsection*{Script — an eye, then a right side of six}
\hlblockfigure{\detokenize{figures/ZH-W471-yan-eye-filmstrip.pdf}}{How \zh{眼} is written, stroke by stroke}

\begin{quote}
\zh{眼}
\end{quote}

\textbf{Two pieces, left before right.} Count the pieces, not the strokes.

\textbf{The left piece, the eye, five strokes.} Like \zh{日}, made tall and narrow, with one bar more: a vertical down the left, a stroke that runs right and turns down, two short bars inside, then the bar across the foot.

\textbf{The right piece, six strokes.} A stroke that runs right and turns down. A bar inside, and a second one under it. Then a vertical down the left that turns up to the right at the foot, a short stroke falling to the left, and last a long stroke falling to the right.

5 + 6 = eleven.

\subsection*{Your turn}
\begin{itemize}
\raggedright
  \item \emph{Write it:} \zh{眼}, all eleven, with the model in front of you
  \item \emph{Cover:} write it again from memory, and count the pieces, not the strokes
\end{itemize}

\begin{tcolorbox}[breakable,colback=teal!4,colframe=teal!35!black,title={Before you move on}]
How many pieces does \zh{眼} have? (\textbf{Two: five strokes and six}.) What is the left piece like? (\textbf{\zh{日}, with one bar more}.)

Source: \href{https://raw.githubusercontent.com/chanind/hanzi-writer-data/68d10a4b21150cae5e1ebbd223eed289cf32d90c/data/\%E7\%9C\%BC.json}{Hanzi Writer Data \zh{眼}, pinned stroke paths and medians}.
\end{tcolorbox}
\section[\texorpdfstring{\zh{小心眼} (xiǎoxīnyǎn)}{小心眼 (xiǎoxīnyǎn)}]{\zh{小心眼} (xiǎoxīnyǎn) — petty; narrow-minded}

\label{lesson:ZH-C471-xiaoxinyan}



\begin{quote}
Before the new one: say the Mandarin for tears, then the Mandarin for right before one's eyes.
\end{quote}

\subsection*{You'll want to know: \zh{小心眼}}
\textbf{\zh{小心眼}} — \emph{xiǎoxīnyǎn} — \textquotedblleft{}petty; narrow-minded\textquotedblright{}.

\begin{quote}
\textbf{\zh{他有点儿小心眼}} \emph{tā yǒudiǎnr xiǎoxīnyǎn} — He's a bit petty.
\end{quote}

\begin{cousinweb}
\textbf{\zh{小心眼}} (\emph{xiǎoxīnyǎn}) is built from \zh{小} 'small', \zh{心} 'heart' and \zh{眼} 'eye', all inherited from Middle Chinese, \emph{sjewX}, \emph{sim} and \emph{ngeanX} (Guangyun, in Baxter's transcription): a heart with a small eye, which sees only small things.
\end{cousinweb}

\begin{grammarlens}[title={What you've built}]
One more word for saying how it was.
\end{grammarlens}

\subsection*{Your turn}
\begin{itemize}
\raggedright
  \item \emph{Say it:} \emph{xiǎoxīnyǎn}
  \item \emph{Say it:} \emph{xiǎoxīnyǎn}, once more
  \item \emph{Say it:} \emph{tā yǒudiǎnr xiǎoxīnyǎn} — He's a bit petty.
  \item \emph{Recall:} say the Mandarin for tears, then the Mandarin for right before one's eyes, then say \emph{xiǎoxīnyǎn} again
\end{itemize}

\begin{tcolorbox}[breakable,colback=teal!4,colframe=teal!35!black,title={Before you move on}]
What does \zh{小心眼} mean? (\textquotedblleft{}Petty; narrow-minded\textquotedblright{}.) Say it once more.
\end{tcolorbox}
\section[jìnshì]{jìnshì — short-sighted}

\label{lesson:ZH-C471-jinshi}



\begin{quote}
Before the new one: say the Mandarin for right before one's eyes, then the Mandarin for petty and narrow-minded.
\end{quote}

\subsection*{You'll want to know: jìnshì}
\emph{jìnshì} — \textquotedblleft{}short-sighted\textquotedblright{}.

\begin{quote}
\emph{wǒ gǎnjué yǎnjing bú tài hǎo kěnéng jìnshì le} — My eyes don't feel right; I might be getting short-sighted.
\end{quote}

Its characters come later; for now it is a word for the ear.

\begin{cousinweb}
\emph{jìnshì} is built from \emph{jìn} 'near' and \emph{shì} 'to look', both inherited from Middle Chinese, \emph{gj+nX} and \emph{dzyijH} (Guangyun, in Baxter's transcription), which glosses the first 'close by' and the second 'to view'.
\end{cousinweb}

\begin{grammarlens}[title={What you've built}]
One more word for saying how it was.
\end{grammarlens}

\subsection*{Your turn}
\begin{itemize}
\raggedright
  \item \emph{Say it:} \emph{jìnshì}
  \item \emph{Say it:} \emph{jìnshì}, once more
  \item \emph{Say it:} \emph{wǒ gǎnjué yǎnjing bú tài hǎo kěnéng jìnshì le} — My eyes don't feel right; I might be getting short-sighted.
  \item \emph{Recall:} say the Mandarin for right before one's eyes, then the Mandarin for petty and narrow-minded, then say \emph{jìnshì} again
\end{itemize}

\begin{tcolorbox}[breakable,colback=teal!4,colframe=teal!35!black,title={Before you move on}]
What does jìnshì mean? (\textquotedblleft{}Short-sighted\textquotedblright{}.) Say it once more.
\end{tcolorbox}
\section[zhǎyǎn]{zhǎyǎn — to blink; to wink}

\label{lesson:ZH-C471-zhayan}



\begin{quote}
Before the new one: say the Mandarin for petty and narrow-minded, then the Mandarin for short-sighted.
\end{quote}

\subsection*{You'll want to know: zhǎyǎn}
\emph{zhǎyǎn} — \textquotedblleft{}to blink; to wink\textquotedblright{}.

\begin{quote}
\emph{tā duì wǒ zhǎyǎn} — He winked at me.
\end{quote}

Its characters come later; for now it is a word for the ear.

Two third tones meet in this word: zhǎ is said with a rising tone before yǎn, and the spelling keeps both marks.

\begin{cousinweb}
\emph{zhǎyǎn} is built from \emph{zhǎ} and \zh{眼} 'eye', both inherited from Middle Chinese, \emph{tsreap} and \emph{ngeanX} (Guangyun, in Baxter's transcription), which glosses the first 'the eye moving'.
\end{cousinweb}

\begin{grammarlens}[title={What you've built}]
One more word for saying how it was.
\end{grammarlens}

\subsection*{Your turn}
\begin{itemize}
\raggedright
  \item \emph{Say it:} \emph{zhǎyǎn}
  \item \emph{Say it:} \emph{zhǎyǎn}, once more
  \item \emph{Say it:} \emph{tā duì wǒ zhǎyǎn} — He winked at me.
  \item \emph{Recall:} say the Mandarin for petty and narrow-minded, then the Mandarin for short-sighted, then say \emph{zhǎyǎn} again
\end{itemize}

\begin{tcolorbox}[breakable,colback=teal!4,colframe=teal!35!black,title={Before you move on}]
What does zhǎyǎn mean? (\textquotedblleft{}To blink; to wink\textquotedblright{}.) Say it once more.
\end{tcolorbox}
